\SourceLeaf{314}%
mit dem Modul der Logarithmen nach Briggs, wobei \(\rho\) in der in Artikel 16 angegebenen Bedeutung zu verstehen ist (und damit auch \(\mu\)). Für diese konstanten Faktoren gilt
\[
\begin{aligned}
\log\mu&=7.9297527989\,(-20)\\
\log\nu&=4.9206912908\,(-30)
\end{aligned}
\]

\indent Bis zu den Gliedern zweiter Ordnung stimmen diese Ergebnisse mit denen aus Artikel 17 überein. Die vorangehende weitere Entwicklung sollte nur deutlich machen, dass die Glieder zweiter Ordnung selbst für die genaueste Rechnung völlig ausreichen. Tatsächlich erreichen die Glieder vierter Ordnung im gesamten hannoverschen Dreiecksnetz in keinem Fall zwei Einheiten der zehnten Dezimalstelle. Nur in ein paar Fällen überschreiten sie \textit{eine} Einheit der zehnten Dezimalstelle.

\Article{20}
Will man mit unseren Formeln, die vom Mittelwert der Breiten und Azimute an beiden Orten ausgehen und nicht von der Breite und dem Azimut an einem Ort, die zu Beginn dieser Abhandlung genannte Aufgabe lösen, so muss man indirekt vorgehen, genauer gesagt durch schrittweise Annäherung, die man beliebig weit treiben kann. Man beginnt mit irgendeinem Näherungswert von \(T\) (fehlt jede andere Kenntnis oder Schätzung, so kann man zunächst das gegebene Azimut am ersten Ort nehmen) und leitet daraus einen viel genaueren Wert ab. Mit diesem wiederholt man dann dieselbe Rechnung und fährt so lange fort, bis sich die Ergebnisse nicht mehr ändern. Bei den ersten Rechnungen braucht man nur 4 oder 5 Stellen der Logarithmen zu berücksichtigen und darf \(\beta\) und \(\tau\) anstelle der korrigierten \(b\) und \(t\) verwenden. Um \(\lambda\) und \(l\) braucht man sich daher bei diesen ersten Rechnungen noch nicht zu kümmern. Bezeichnet man am ersten Ort die Breite mit \(B^0\) und das Azimut mit \(T^0\), so lauten die Formeln der Reihe nach:
\[
\begin{aligned}
\beta&=r\cos T\\
B&=B^0-\tfrac12\beta\\
\tau&=r\sin T\operatorname{tang}B\\
T&=T^0-\tfrac12\tau
\end{aligned}
\]
